% =====================================================================
%  exam-frequency-report / assets/report-preamble.tex
%  Reusable preamble for an exam-frequency report.
%
%  Compile with XeLaTeX (twice).  Load-bearing choices, each of which
%  fixes a real failure mode -- see the skill's Pitfalls table:
%    * ctexart + xelatex          -> CJK actually renders
%    * amsmath only, before newtxmath -> avoids the \Bbbk clash with amssymb
%    * xcolor[table]              -> \rowcolor is defined
%    * \noindent before \rule/\tabularx -> no 17pt \parindent overfull
%  Adjust \setReportMeta and the geometry block per task.
% =====================================================================
\documentclass[11pt]{article}

% --- page geometry ---------------------------------------------------
% Keep the long edge at a size you have already validated (A4 height
% preserves the vertical layout).  To match a reading device, set
%   paperwidth = paperheight * Wpx / Hpx
% e.g. 13-inch iPad Air M2 portrait (2048x2732):
%   paperheight=297mm, paperwidth=297*512/683=222.6413mm
\usepackage[paperwidth=222.6413mm,paperheight=297mm,
            left=1.9cm,right=1.9cm,top=2.1cm,bottom=1.7cm,headsep=0.5cm]{geometry}

% --- math: order matters ---------------------------------------------
% DO NOT load newtxtext in a CJK (ctex) document.
% On macOS ctex falls back to the macold fontset, whose \ttfamily CJK slot
% points at STFangsong -- absent from current macOS.  Any Chinese inside
% \texttt (e.g. a path written as \code{试卷/教材}) then fails with
%   ! Font TU/STFangsong ... not loadable
% and those glyphs come out as mojibake.  Verified: newtxtext triggers it,
% newtxmath is harmless, and \setCJKmonofont does NOT repair it.
% Pure-ASCII (non-ctex) documents may load newtxtext normally.
\usepackage{amsmath}
\usepackage{newtxmath}

% --- tables, colour, boxes -------------------------------------------
\usepackage[table]{xcolor}
\usepackage{booktabs}
\usepackage{longtable}
\usepackage{tabularx}
\usepackage{array}
\usepackage{multirow}
\usepackage{makecell}
\usepackage{ragged2e}
\usepackage{enumitem}
\usepackage{tcolorbox}
\tcbuselibrary{skins,breakable}
\usepackage{pgfplots}
\pgfplotsset{compat=1.18}
\usetikzlibrary{patterns}
\usepackage{titlesec}
\usepackage{fancyhdr}
\usepackage{caption}
\usepackage[hidelinks,bookmarksnumbered]{hyperref}

% --- palette ---------------------------------------------------------
\definecolor{ink}{HTML}{1F2933}
\definecolor{deepblue}{HTML}{1B3A5C}
\definecolor{accentA}{HTML}{B23A48}   % S tier
\definecolor{accentB}{HTML}{D97706}   % A tier
\definecolor{accentC}{HTML}{2F6F4E}   % B tier
\definecolor{accentD}{HTML}{6B7280}   % C tier
\definecolor{tintA}{HTML}{FBE9EB}
\definecolor{tintB}{HTML}{FDF2E3}
\definecolor{tintC}{HTML}{E8F2EC}
\definecolor{tintD}{HTML}{F0F1F3}
\definecolor{rulegray}{HTML}{C9CED6}
\definecolor{tabhead}{HTML}{1B3A5C}

% --- headings --------------------------------------------------------
\titleformat{\section}{\normalfont\Large\bfseries\color{deepblue}}{\thesection}{0.7em}{}
\titleformat{\subsection}{\normalfont\large\bfseries\color{deepblue!85!black}}{\thesubsection}{0.6em}{}
\titlespacing*{\section}{0pt}{1.6em}{0.7em}
\titlespacing*{\subsection}{0pt}{1.2em}{0.5em}

% --- running head ----------------------------------------------------
\pagestyle{fancy}
\fancyhf{}
\renewcommand{\headrulewidth}{0.4pt}
\fancyhead[L]{\footnotesize\color{deepblue}\ReportShortTitle}
\fancyhead[R]{\footnotesize\color{deepblue}\thepage}
\fancyfoot[C]{\scriptsize\color{accentD}\ReportSourceLine}
\newcommand{\ReportShortTitle}{Exam-frequency report}
\newcommand{\ReportSourceLine}{}

% --- helpers ---------------------------------------------------------
% Section-code tag.  Use the textbook's OWN numbering.
\newcommand{\secode}[1]{\texttt{\color{deepblue}#1}}
\newcommand{\tierS}{\textbf{\color{accentA}S}}
\newcommand{\tierA}{\textbf{\color{accentB}A}}
\newcommand{\tierB}{\textbf{\color{accentC}B}}
\newcommand{\tierC}{\textbf{\color{accentD}C}}
\newcommand{\hot}[1]{\textcolor{accentA}{\bfseries #1}}
\newcommand{\code}[1]{\texttt{\small #1}}
\newcommand{\thd}[1]{\textcolor{white}{\bfseries #1}}
\newcolumntype{L}[1]{>{\RaggedRight\arraybackslash}p{#1}}
\newcolumntype{C}[1]{>{\Centering\arraybackslash}p{#1}}
% NOTE: \newlist needs an explicit label= or enumitem raises "Undefined label".
\newlist{qparts}{enumerate}{1}
\setlist[qparts]{label=(\alph*),leftmargin=1.5em,itemsep=2.5pt,topsep=3pt,parsep=0pt}
\newlist{qnum}{enumerate}{1}
\setlist[qnum]{label=(\arabic*),leftmargin=1.5em,itemsep=2.5pt,topsep=3pt,parsep=0pt}
\newlist{bul}{itemize}{1}
\setlist[bul]{label=\textbullet,leftmargin=1.6em,itemsep=1.5pt,topsep=2pt,parsep=0pt}
\captionsetup{font={small,color=deepblue},labelfont=bf,skip=5pt}

% --- callout boxes ---------------------------------------------------
\newtcolorbox{keybox}[1][]{colback=tintC,colframe=accentC,boxrule=0.7pt,arc=2pt,
  left=8pt,right=8pt,top=6pt,bottom=6pt,fonttitle=\bfseries\color{white},
  coltitle=white,colbacktitle=accentC,breakable,#1}
\newtcolorbox{warnbox}[1][]{colback=tintA,colframe=accentA,boxrule=0.7pt,arc=2pt,
  left=8pt,right=8pt,top=6pt,bottom=6pt,fonttitle=\bfseries\color{white},
  coltitle=white,colbacktitle=accentA,breakable,#1}
\newtcolorbox{infobox}[1][]{colback=tintB,colframe=accentB,boxrule=0.7pt,arc=2pt,
  left=8pt,right=8pt,top=6pt,bottom=6pt,fonttitle=\bfseries\color{white},
  coltitle=white,colbacktitle=accentB,breakable,#1}

% --- one-item-per-page layout (review sets) --------------------------
% tcolorbox `height fill` absorbs the remaining page height.
% Do NOT combine with `breakable`.
\newtcolorbox{workspace}{height fill,enhanced,colback=white,colframe=gray!45,
  boxrule=0.5pt,arc=1.5pt,left=9pt,right=9pt,top=6pt,bottom=6pt}

% --- pgfplots reference lines on a symbolic axis ---------------------
% A numeric line on `symbolic x coords` fails with
%   "the input coordinate '...' has not been defined".
% Draw it in axis-description coordinates instead, e.g. y=33.75 of ymax=60:
%   \draw[accentA,dashed,thick]
%     (axis description cs:0,0.5625) -- (axis description cs:1,0.5625);
%   \node[anchor=south west,font=\scriptsize,color=accentA]
%     at (axis description cs:0.02,0.578) {four-year mean};


\setlength{\parindent}{0pt}\setlength{\parskip}{4pt}\setlength{\headheight}{14pt}
\setlist[enumerate]{label=(\arabic*),leftmargin=1.8em,itemsep=4pt}
\newcommand{\E}{\mathbb E}\newcommand{\Var}{\operatorname{Var}}\newcommand{\Cov}{\operatorname{Cov}}
\newcommand{\key}[1]{\par\vspace{3pt}\textbf{\color{deepblue}#1}\par}
\newcommand{\qtitle}[3]{\clearpage
\begin{tcolorbox}[colback=deepblue,colframe=deepblue,coltext=white,boxrule=0pt,arc=2pt]
{\Large\bfseries #1}\hfill {Day #2}\end{tcolorbox}
{\large\bfseries #3}\par\vspace{2mm}}

\renewcommand{\ReportShortTitle}{Chapter 2 Review Problem Set}\renewcommand{\ReportSourceLine}{SSA · Chapter 2 · 21 problems · 2026-10-13}\begin{document}\thispagestyle{empty}\vspace*{15mm}
{\Huge\bfseries\color{deepblue}Chapter 2 Review Problem Set\par}\vspace{8mm}
{\Large 21 problems · 7 days · 3 problems per day\par}\vspace{10mm}
\begin{keybox}[title=Scope]
All ten Homework 2 exercises are included. Additional questions target Gaussian moments, correlation identities, finite-dimensional independence, and assumptions that matter in past papers.
Each problem has its own page and handwriting space. A matched Chinese solution key uses the same numbering.
\end{keybox}
\begin{infobox}[title=Source labels]
Textbook problems from Homework 2 retain the existing English translation of the Chinese textbook. This is not an official translation.
Questions labelled ``Adapted from'' restate or extend selected exam tasks; extra assumptions are explicit. Recall questions are newly written.
\end{infobox}
\vspace{8mm}Required homework: 2.3, 2.13, 2.16, 2.19.\\
Optional homework: 2.1, 2.4, 2.7, 2.8, 2.11, 2.18.\\
Homework issued: 13 October 2026. Due: 27 October 2026.\\
This seven-day revision schedule does not change the homework deadline.
\vfill Stochastic Signal Analysis · Chapter 2: Random Signals\\
Based on the fifth-edition Chinese textbook and four local past papers.
\clearpage\section*{Seven-day plan and problem index}


\small\begin{tabularx}{\textwidth}{@{}c l X@{}}\toprule Day & Problems & Focus\\\midrule
1&1--3&Gaussian tools: Gaussian checklist; 2.19; 2.17\\
2&4--6&Gaussian applications: Square-law device; 2.16; 2.18\\
3&7--9&Phase and cross-correlation: Uniform phase; sampled sin/cos; 2.11\\
4&10--12&Correlation and independence: Identities; random line; signal independence\\
5&13--15&Finite-dimensional laws: 2.3; coin-toss exam; 2.1\\
6&16--18&Sequences and linear combinations: 2.4; 2.8; 2.7\\
7&19--21&Operations and assumptions: 2.13; 2.14; white-noise caveat\\
\bottomrule\end{tabularx}\vspace{5mm}\par\begin{tabularx}{\textwidth}{@{}r L{7cm} X@{}}\toprule No. & Source & Textbook section\\\midrule
1 & Recall / extension & 2.4.1--2.4.3\\
2 & Textbook 2.19 / Required & 2.4.1\\
3 & Textbook 2.17 / translated & 2.4.1 / 2.4.3\\
4 & Adapted from 2022--23 Q3 & 2.4.3\\
5 & Textbook 2.16 / Required & 2.4.1 / 2.4.3\\
6 & Textbook 2.18 / Optional & 2.4.1 / 2.4.3\\
7 & Adapted from 2022--23 Q2(a,b) & 2.2.1 / 2.5\\
8 & Adapted from 2024--25 Q1 & 2.3.2 / 2.5\\
9 & Textbook 2.11 / Optional & 2.3.2 / 2.5\\
10 & Recall / extension & 2.1.3 / 2.3.3\\
11 & Adapted from 2022--23 Q1 & 2.1.1--2.1.3\\
12 & Recall / extension & 2.3.1 / 2.5\\
13 & Textbook 2.3 / Required & 2.1.2 / 2.3.1\\
14 & Adapted from 2025--26 Q1 & 2.1.2 / 2.3.1 / 2.5\\
15 & Textbook 2.1 / Optional & 2.1.1 / 2.2.1\\
16 & Textbook 2.4 / Optional & 2.2.2\\
17 & Textbook 2.8 / Optional & 2.2.2 / 2.1.3\\
18 & Textbook 2.7 / Optional & 2.4.2 / 2.1.3\\
19 & Textbook 2.13 / Required & 2.2.1 / 2.3.4\\
20 & Textbook 2.14 / translated & 2.3.4\\
21 & Adapted from 2023--24 Q2(c) & 2.4.3\\
\bottomrule\end{tabularx}\normalsize
\qtitle{Problem 1}{1}{Gaussian checklist}
{\small\color{accentA}(Recall / extension)}\par{\small\color{accentD}Textbook section: 2.4.1--2.4.3}\par\vspace{2mm}
Define a Gaussian process using its finite-dimensional distributions. For a Gaussian vector with mean $\boldsymbol\mu$ and covariance $C$, state its characteristic function, its ordinary PDF (including the condition for its existence), and the law of $G\boldsymbol X+\boldsymbol b$.

Is marginal Gaussianity sufficient for joint Gaussianity? Give a counterexample. State when uncorrelatedness implies independence.
\par\vspace{3mm}\begin{workspace}\textcolor{gray}{Working / Answer}\end{workspace}
\qtitle{Problem 2}{1}{Joint Gaussian PDF and CF}
{\small\color{accentA}(Textbook 2.19 / Required)}\par{\small\color{accentD}Textbook section: 2.4.1}\par\vspace{2mm}
Let the random vector $(X,Y)\sim N(\boldsymbol{\mu},\boldsymbol{C})$, where
\[
\boldsymbol{\mu}=\begin{bmatrix}2\\2\end{bmatrix},\qquad
\boldsymbol{C}=\begin{bmatrix}2&3\\3&5\end{bmatrix}.
\]
Find the joint probability density function of $(X,Y)$ and its characteristic function $\phi_{XY}(u,v)$.
\par\vspace{3mm}\begin{workspace}\textcolor{gray}{Working / Answer}\end{workspace}
\qtitle{Problem 3}{1}{Fourth-order Gaussian moment}
{\small\color{accentA}(Textbook 2.17 / translated)}\par{\small\color{accentD}Textbook section: 2.4.1 / 2.4.3}\par\vspace{2mm}
Let $X(t)$ be a zero-mean Gaussian signal and put $X_k=X(t_k)$ for $k=1,2,3,4$. Prove that
\[\begin{aligned}E[X_1X_2X_3X_4]={}&E[X_1X_2]E[X_3X_4]\\&+E[X_1X_3]E[X_2X_4]+E[X_1X_4]E[X_2X_3].\end{aligned}\]
\par\vspace{3mm}\begin{workspace}\textcolor{gray}{Working / Answer}\end{workspace}
\qtitle{Problem 4}{2}{Gaussian signal through a square-law device}
{\small\color{accentA}(Adapted from 2022--23 Q3)}\par{\small\color{accentD}Textbook section: 2.4.3}\par\vspace{2mm}
A zero-mean stationary Gaussian white noise $N(t)$ has PSD $N_0/2$, where $N_0>0$. It passes through an RC low-pass filter
\[H(j\omega)=\frac1{1+j\omega RC},\qquad R,C>0,\]
whose output is $X(t)$. A second device produces $Y(t)=aX^2(t)$, where $a$ is real.
\begin{enumerate}\item Find $R_X(\tau)$.\item Find $E[Y(t)]$, $R_Y(\tau)$ and $C_Y(\tau)$.\item Is $Y(t)$ Gaussian when $a\ne0$?\end{enumerate}
The transfer function is supplied in place of the original circuit figure; the covariance and Gaussianity questions are extensions.
\par\vspace{3mm}\begin{workspace}\textcolor{gray}{Working / Answer}\end{workspace}
\qtitle{Problem 5}{2}{Gaussian sinusoid}
{\small\color{accentA}(Textbook 2.16 / Required)}\par{\small\color{accentD}Textbook section: 2.4.1 / 2.4.3}\par\vspace{2mm}
Let the random process be
\[
X(t)=A\cos\omega t+B\sin\omega t,
\]
where $\omega$ is a constant and $A$ and $B$ are independent Gaussian random variables. Given that
\[
E[A]=E[B]=0,\qquad E[A^2]=E[B^2]=\sigma^2,
\]
find the one-dimensional and two-dimensional probability density functions of $X(t)$.
\par\vspace{3mm}\begin{workspace}\textcolor{gray}{Working / Answer}\end{workspace}
\qtitle{Problem 6}{2}{Singular three-dimensional Gaussian law}
{\small\color{accentA}(Textbook 2.18 / Optional)}\par{\small\color{accentD}Textbook section: 2.4.1 / 2.4.3}\par\vspace{2mm}
A Gaussian signal has mean and covariance function
\[
m_X(t)=2,\qquad C_X(t_1,t_2)=8\cos(t_1-t_2).
\]
Write down its three-dimensional probability density function at the times
\[
t_1=0,\qquad t_2=0.5,\qquad t_3=1.
\]
% Preserve the textbook question. A later solution must address the singular covariance matrix.
\par\vspace{3mm}\begin{workspace}\textcolor{gray}{Working / Answer}\end{workspace}
\qtitle{Problem 7}{3}{Uniform random phase}
{\small\color{accentA}(Adapted from 2022--23 Q2(a,b))}\par{\small\color{accentD}Textbook section: 2.2.1 / 2.5}\par\vspace{2mm}
Let $\eta(t)=\xi(t)\cos(\omega_0t+\Theta)$, where $\xi(t)$ is a real WSS process with finite second moment and $\Theta\sim U[0,2\pi]$ is independent of the entire process.
Find the mean and autocorrelation of $\eta(t)$, and determine whether it is WSS. Show the phase averages explicitly. No PSD calculation is required.
\par\vspace{3mm}\begin{workspace}\textcolor{gray}{Working / Answer}\end{workspace}
\qtitle{Problem 8}{3}{Orthogonal but dependent signals}
{\small\color{accentA}(Adapted from 2024--25 Q1)}\par{\small\color{accentD}Textbook section: 2.3.2 / 2.5}\par\vspace{2mm}
Let
\[X[n]=\sin(2\pi n+\Phi),\qquad Y[n]=\cos(2\pi n+\Phi),\qquad n\in\mathbb Z,\]
where $\Phi\sim U(0,\pi)$.
\begin{enumerate}\item Find $R_{XY}(m,n)$ and the two means.\item Are the signals orthogonal, uncorrelated, or statistically independent? Justify each answer.\item Are they jointly stationary?\end{enumerate}
\par\vspace{3mm}\begin{workspace}\textcolor{gray}{Working / Answer}\end{workspace}
\qtitle{Problem 9}{3}{Pairwise versus mutual independence}
{\small\color{accentA}(Textbook 2.11 / Optional)}\par{\small\color{accentD}Textbook section: 2.3.2 / 2.5}\par\vspace{2mm}
Consider the two random signals
\[
X(t)=A\sin(\omega t+\Theta),\qquad
Y(t)=B\cos(\omega t+\Theta),
\]
where $A$ and $B$ are random variables with unspecified distributions, $\Theta$ is uniformly distributed on $[0,2\pi]$, and $\omega$ is a constant. The random variables $A$, $B$, and $\Theta$ are pairwise statistically independent.
% Preserve the textbook's pairwise independence assumption; do not strengthen it to mutual independence.
\begin{enumerate}
\item Find the cross-correlation function $R_{XY}(t_1,t_2)$ of the two random signals.
\item Discuss their orthogonality, uncorrelatedness, and statistical independence.
\end{enumerate}
\par\vspace{3mm}\begin{workspace}\textcolor{gray}{Working / Answer}\end{workspace}
\qtitle{Problem 10}{4}{Correlation identities}
{\small\color{accentA}(Recall / extension)}\par{\small\color{accentD}Textbook section: 2.1.3 / 2.3.3}\par\vspace{2mm}
For a real process $X(t)$ with finite second moments:
\begin{enumerate}\item Define $m_X$, $R_X$ and $C_X$, and relate $R_X(t,t)$ to variance.\item Prove that every sampled covariance matrix is positive semidefinite, and that $|C_X(s,t)|^2\le C_X(s,s)C_X(t,t)$.\item For $Y(t)=X(t+a)-X(t)$, expand $R_Y(s,t)$ without assuming stationarity.\end{enumerate}
\par\vspace{3mm}\begin{workspace}\textcolor{gray}{Working / Answer}\end{workspace}
\qtitle{Problem 11}{4}{A random straight line}
{\small\color{accentA}(Adapted from 2022--23 Q1)}\par{\small\color{accentD}Textbook section: 2.1.1--2.1.3}\par\vspace{2mm}
Let $X(t)=U+t$ for $-2\le t\le2$, where $U\sim U[0,\pi]$.
Sketch two sample functions, find the characteristic function of $X(t)$, and compute its mean and autocovariance. Explain whether samples at distinct times are independent.
\par\vspace{3mm}\begin{workspace}\textcolor{gray}{Working / Answer}\end{workspace}
\qtitle{Problem 12}{4}{Independence at the process level}
{\small\color{accentA}(Recall / extension)}\par{\small\color{accentD}Textbook section: 2.3.1 / 2.5}\par\vspace{2mm}
State the finite-dimensional definition of independence of two random signals. Does independence of $X(t)$ and $Y(t)$ at every equal time imply independence of the two signals?
Use an iid sequence $\{Z_n\}$ of fair signs and $X_n=Z_n$, $Y_n=Z_{n-1}$ to answer. Also state the criterion for two jointly Gaussian signals to be independent.
\par\vspace{3mm}\begin{workspace}\textcolor{gray}{Working / Answer}\end{workspace}
\qtitle{Problem 13}{5}{One toss, two observation times}
{\small\color{accentA}(Textbook 2.3 / Required)}\par{\small\color{accentD}Textbook section: 2.1.2 / 2.3.1}\par\vspace{2mm}
A coin is tossed, and a random process is defined by
\[
X(t)=\begin{cases}
\cos\pi t,&\text{if heads occurs},\\
2t,&\text{if tails occurs}.
\end{cases}
\]
Heads and tails are equally likely. Find:
\begin{enumerate}
\item The one-dimensional distribution functions $F_X(x,1/2)$ and $F_X(x,1)$ of $X(t)$;
\item The two-dimensional distribution function $F_X(x_1,x_2;1/2,1)$ of $X(t)$;
\item Sketch the graphs of the above distribution functions.
\end{enumerate}
\par\vspace{3mm}\begin{workspace}\textcolor{gray}{Working / Answer}\end{workspace}
\qtitle{Problem 14}{5}{One toss versus fresh tosses}
{\small\color{accentA}(Adapted from 2025--26 Q1)}\par{\small\color{accentD}Textbook section: 2.1.2 / 2.3.1 / 2.5}\par\vspace{2mm}
A fair coin is tossed once. Define $X(t)=\sin(2\pi t)$ for heads and $X(t)=t^2$ for tails.
\begin{enumerate}\item Find the generalized PDF of $X(2)$.\item Find the joint generalized PDF of $(X(1),X(2))$.\item Find $C_X(s,t)$.\item Now define $Y[n]=\sin(2\pi n)$ or $n^2$ by a fresh independent fair toss at each integer $n$. Find the joint generalized PDF of $(Y[1],Y[2])$.\end{enumerate}
\par\vspace{3mm}\begin{workspace}\textcolor{gray}{Working / Answer}\end{workspace}
\qtitle{Problem 15}{5}{Random amplitude and a degenerate time}
{\small\color{accentA}(Textbook 2.1 / Optional)}\par{\small\color{accentD}Textbook section: 2.1.1 / 2.2.1}\par\vspace{2mm}
Consider the sinusoidal random signal $X(t)=V\cos\omega t$, where $0\le t<\infty$, $\omega$ is a constant, and $V$ is uniformly distributed on $[0,1]$.
\begin{enumerate}
\item Sketch two sample functions of this process.
\item Determine the probability density function of the random variable $X(t_i)$ at each of the times
\[
t_i=0,\quad\frac{\pi}{4\omega},\quad\frac{3\pi}{4\omega},\quad\frac{\pi}{\omega}.
\]
\item Find the probability density function of $X(t_0)$ when $t_0=\pi/(2\omega)$.
\end{enumerate}
\par\vspace{3mm}\begin{workspace}\textcolor{gray}{Working / Answer}\end{workspace}
\qtitle{Problem 16}{6}{Bernoulli words}
{\small\color{accentA}(Textbook 2.4 / Optional)}\par{\small\color{accentD}Textbook section: 2.2.2}\par\vspace{2mm}
Suppose that the binary data sequence $\{B(n),n=1,2,3,\ldots\}$ is a Bernoulli random sequence. Each bit corresponds to a random variable $B(n)$, with
\[
P[B(n)=0]=0.2,\qquad P[B(n)=1]=0.8.
\]
\begin{enumerate}
\item What is the probability that four consecutive bits form the word $\{1011\}$?
\item What is the mean word formed by four consecutive bits?
\item Which word formed by four consecutive bits has the highest probability?
\end{enumerate}
\par\vspace{3mm}\begin{workspace}\textcolor{gray}{Working / Answer}\end{workspace}
\qtitle{Problem 17}{6}{A fair-sign sequence}
{\small\color{accentA}(Textbook 2.8 / Optional)}\par{\small\color{accentD}Textbook section: 2.2.2 / 2.1.3}\par\vspace{2mm}
Suppose that the Bernoulli sequence $\{I_n,n=1,2,\ldots\}$ takes the values $-1$ and $+1$ with equal probabilities.

Find its one-dimensional probability density function, mean, and covariance function.
\par\vspace{3mm}\begin{workspace}\textcolor{gray}{Working / Answer}\end{workspace}
\qtitle{Problem 18}{6}{Gaussian straight-line process}
{\small\color{accentA}(Textbook 2.7 / Optional)}\par{\small\color{accentD}Textbook section: 2.4.2 / 2.1.3}\par\vspace{2mm}
The position of a moving particle is described by the straight-line process
\[
X(t)=Vt+X_0,
\]
where $V\sim N(1,1)$ and $X_0\sim N(0,2)$ are independent.

Find the one-dimensional probability density function, mean, and variance of the random variable $X(t)$ at time $t$.
\par\vspace{3mm}\begin{workspace}\textcolor{gray}{Working / Answer}\end{workspace}
\qtitle{Problem 19}{7}{Random voltage across a parallel RC circuit}
{\small\color{accentA}(Textbook 2.13 / Required)}\par{\small\color{accentD}Textbook section: 2.2.1 / 2.3.4}\par\vspace{2mm}
Consider the sinusoidal voltage signal
\[
X(t)=A\cos(\omega t+\Theta),
\]
where $\omega$ is a constant, $A\sim U(-1,+1)$, and $\Theta\sim U(-\pi,+\pi)$. The random variables $A$ and $\Theta$ are independent.

If this voltage is applied across a parallel RC circuit, find the total current signal and its mean-square value.
\par\vspace{3mm}\begin{workspace}\textcolor{gray}{Working / Answer}\end{workspace}
\qtitle{Problem 20}{7}{Integration of a random signal}
{\small\color{accentA}(Textbook 2.14 / translated)}\par{\small\color{accentD}Textbook section: 2.3.4}\par\vspace{2mm}
For the integral of a random signal,
\[Y(t)=\int_0^t X(u)\,du,\qquad t\ge0,\]
find expressions for its mean and autocorrelation function. State sufficient conditions that justify interchanging integration and expectation.
\par\vspace{3mm}\begin{workspace}\textcolor{gray}{Working / Answer}\end{workspace}
\qtitle{Problem 21}{7}{White noise is not a Gaussian assumption}
{\small\color{accentA}(Adapted from 2023--24 Q2(c))}\par{\small\color{accentD}Textbook section: 2.4.3}\par\vspace{2mm}
A zero-mean white noise $W(t)$ with PSD $N_0/2$ passes through $h(t)=e^{-t}u(t)$ to produce $X(t)$. Let $Y(t)=X(t)-X(t-T)$, with $T>0$ and $N_0>0$.
You may use $R_X(\tau)=N_0e^{-|\tau|}/4$.
\begin{enumerate}\item Do the stated assumptions imply that $Y(t)$ is Gaussian?\item Under the additional assumption that $W$ is Gaussian white noise, find the one-dimensional PDF of $Y(t)$.\item Describe the case $T=0$.\end{enumerate}
The extra Gaussian assumption in (2) is explicit; it was not stated in the original exam stem.
\par\vspace{3mm}\begin{workspace}\textcolor{gray}{Working / Answer}\end{workspace}
\end{document}